% =====================================================================
\section{Seminario de problemas}
\label{sec:seminario}
% =====================================================================

Los once problemas de la hoja se agrupan en cinco tipos.

\begin{table}[ht]
\centering
\caption{Mapa de métodos del seminario.}
\label{tab:mapa}
\small
\begin{tabular}{@{}l l l@{}}
\toprule
\color{cian}Tipo & \color{cian}Problemas & \color{cian}Herramienta clave\\
\midrule
A. Centro de masas & 1, 3, 4 & \eqref{eq:cm}, simetría, elemento $\dd m$\\
B. Momentos y König & 2 & \cref{teo:konig1}\\
C. Conservación de $\vb{P}$ y movimiento del CM & 5, 6, 7 & \cref{teo:momento-lineal}, energía\\
D. Masa variable & 8, 9 & \cref{teo:masa-variable}\\
E. Virial y Pappus & 10, 11 & \cref{cor:virial-potencia}, \cref{teo:pappus2}\\
\bottomrule
\end{tabular}
\end{table}

% =====================================================================
\subsection{Tipo A · Centro de masas}
% =====================================================================

\begin{problema}{1}[Triángulo equilátero]
\label{prob:1}
Calcular el centro de masas de un sistema de tres partículas de masa $m$
situadas en los vértices de un triángulo equilátero de lado $L$.
\end{problema}

\begin{resolucion}
\begin{enumerate}
  \item \emph{Coordenadas} (\cref{fig:triangulo}):
  $\vb{r}_{1}=(0,0)$, $\vb{r}_{2}=(L,0)$, $\vb{r}_{3}=\bigl(\tfrac L2,\tfrac{\sqrt3}{2}L\bigr)$; $M=3m$.
  \item \emph{Definición} \eqref{eq:cm}:
  \[
    \Rc=\frac{m}{3m}\Bigl[(0,0)+(L,0)+\bigl(\tfrac L2,\tfrac{\sqrt3}{2}L\bigr)\Bigr]
    =\resultado{\Bigl(\frac L2,\ \frac{\sqrt3}{6}L\Bigr)} .
  \]
  \item \emph{Interpretación}: con masas iguales, el CM es la media de las
  posiciones, es decir, el baricentro del triángulo. Está sobre las tres
  medianas (ejes de simetría), a $\tfrac23$ de cada vértice, y su distancia a cada
  vértice es $\abs{\Rc-\vb{r}_{1}}=L\sqrt{\tfrac14+\tfrac1{12}}=L/\sqrt3$.\qedhere
\end{enumerate}
\end{resolucion}

% ---------------------------------------------------------------------
\begin{figure}[ht]
\centering
\begin{tikzpicture}[scale=1.6]
  \coordinate (A) at (0,0); \coordinate (B) at (3,0); \coordinate (C) at (1.5,{1.5*sqrt(3)});
  \coordinate (G) at (1.5,{0.5*sqrt(3)});
  \fill[cian,opacity=0.07] (A) -- (B) -- (C) -- cycle;
  \draw[texto!80,line width=0.8pt] (A) -- (B) -- (C) -- cycle;
  \draw[aux] (A) -- ($(B)!0.5!(C)$);
  \draw[aux] (B) -- ($(A)!0.5!(C)$);
  \draw[aux] (C) -- ($(A)!0.5!(B)$);
  \foreach \P/\l/\a in {A/1/north east,B/2/north west,C/3/south}{
    \node[punto,inner sep=1.8pt] at (\P) {};
    \node[etiq,anchor=\a,inner sep=3pt] at (\P) {$m$};}
  \node[punto,fill=violeta,inner sep=1.8pt] at (G) {};
  \node[violeta,font=\small,anchor=south west,inner sep=1pt] at ($(G)+(60:0.22)$) {CM};
  \draw[violeta,line width=0.5pt,<->] (3.35,0) -- (3.35,{0.5*sqrt(3)});
  \draw[aux] (G) -- (3.45,{0.5*sqrt(3)});
  \node[violeta,font=\small,anchor=west,inner sep=3pt] at (3.35,{0.25*sqrt(3)}) {$\frac{\sqrt3}{6}L$};
  \node[peq,anchor=north,inner sep=4pt] at (1.5,0) {$L$};
\end{tikzpicture}
\caption{Tres masas iguales en un triángulo equilátero: el CM es el
baricentro, intersección de las medianas, a altura $\sqrt3L/6$ sobre la base.}
\label{fig:triangulo}
\end{figure}
% ---------------------------------------------------------------------

\begin{problema}{3}[Varilla de densidad lineal variable]
\label{prob:3}
Determinar el centro de masas de una varilla rectilínea de longitud $L$ y de
sección transversal $S$ con densidad $\lambda=\lambda_{0}+kx$.
\end{problema}

\begin{resolucion}
Se toma la varilla sobre $x\in[0,L]$ y $\lambda$ como densidad \emph{lineal}
(masa por unidad de longitud), de modo que $\dd m=\lambda\,\dd x$.
\begin{observacion}[¿Y la sección $S$?]
Si la densidad dada fuera volúmica, $\rho=\rho_{0}+kx$, sería $\dd m=\rho S\,\dd x$, es
decir, $\lambda=\rho S$. Como $S$ es constante, aparece en $M$ y en $\int x\,\dd m$ y se
cancela en $X$: el resultado es el mismo.
\end{observacion}
\begin{enumerate}
  \item \emph{Simetría}: el CM está sobre el eje de la varilla. Como $\lambda$
  crece con $x$, estará desplazado hacia $x=L$.
  \item \emph{Masa}:
  $M=\displaystyle\int_{0}^{L}(\lambda_{0}+kx)\,\dd x=\lambda_{0}L+\frac{kL^{2}}{2}$.
  \item \emph{Primer momento}:
  $\displaystyle\int x\,\dd m=\int_{0}^{L}(\lambda_{0}x+kx^{2})\,\dd x=\frac{\lambda_{0}L^{2}}{2}+\frac{kL^{3}}{3}$.
  \item \emph{Cociente}:
  \[
    X=\frac{\tfrac{\lambda_{0}L^{2}}{2}+\tfrac{kL^{3}}{3}}{\lambda_{0}L+\tfrac{kL^{2}}{2}}
    =\resultado{\frac{L}{3}\,\frac{3\lambda_{0}+2kL}{2\lambda_{0}+kL}} .
  \]
  \item \emph{Comprobación}: $k=0$ da $X=L/2$ (varilla homogénea); $\lambda_{0}=0$
  da $X=2L/3$ (densidad $\propto x$, como un triángulo); y si $k>0$, $X\in(L/2,2L/3)$.\qedhere
\end{enumerate}
\end{resolucion}

\begin{problema}{4}[Cascarón hemisférico]
\label{prob:4}
Determinar el centro de masas de un cascarón hemisférico de radio $R$ y
densidad superficial constante.
\end{problema}

\begin{resolucion}
\begin{enumerate}
  \item \emph{Simetría}: el eje $z$ es de revolución, así que $X=Y=0$.
  \item \emph{Elemento} (esféricas con $r=R$, \cref{fig:cascaron}):
  $\dd S=R^{2}\sin\theta\,\dd\theta\,\dd\varphi$, $\dd m=\sigma\,\dd S$, $z=R\cos\theta$,
  con $\theta\in[0,\tfrac\pi2]$ y $\varphi\in[0,2\pi)$.
  \item \emph{Masa}:
  $M=\sigma R^{2}\displaystyle\int_{0}^{2\pi}\!\!\dd\varphi\int_{0}^{\pi/2}\!\!\sin\theta\,\dd\theta=2\pi\sigma R^{2}$.
  \item \emph{Primer momento}:
  \[
    \int z\,\dd m=\sigma R^{3}\int_{0}^{2\pi}\!\!\dd\varphi\int_{0}^{\pi/2}\!\!\cos\theta\sin\theta\,\dd\theta
    =2\pi\sigma R^{3}\Bigl[\tfrac12\sin^{2}\theta\Bigr]_{0}^{\pi/2}=\pi\sigma R^{3}.
  \]
  \item \emph{Resultado}: $Z=\dfrac{\pi\sigma R^{3}}{2\pi\sigma R^{2}}=\resultado{\dfrac{R}{2}}$.
  \item \emph{Comparación}: la semiesfera \emph{maciza} tiene $Z=3R/8<R/2$: en el
  cascarón toda la masa está lejos del centro.\qedhere
\end{enumerate}
\end{resolucion}

% ---------------------------------------------------------------------
\begin{figure}[ht]
\centering
\begin{tikzpicture}[scale=1.5]
  \def\R{1.8}\def\th{48}
  \fill[cian,opacity=0.07] (\R,0) arc[start angle=0,end angle=180,radius=\R] -- cycle;
  \draw[texto!85,line width=1.1pt] (\R,0) arc[start angle=0,end angle=180,radius=\R];
  \draw[texto!60,line width=0.6pt] (0,0) ellipse[x radius=\R,y radius=0.32];
  \draw[eje] (0,0) -- (0,2.35) node[anchor=south,peq,inner sep=2pt]{$z$};
  % banda dS
  \draw[cian,line width=3pt] ({\R*cos(90-\th)},{\R*sin(90-\th)}) arc[start angle={90-\th},end angle={90-\th+8},radius=\R];
  \draw[aux] (0,0) -- ({\R*cos(90-\th)},{\R*sin(90-\th)});
  \draw[violeta,line width=0.6pt,->] (0,0.5) arc[start angle=90,end angle={90-\th},radius=0.5];
  \node[violeta,font=\small] at ({0.68*cos(90-\th/2)},{0.68*sin(90-\th/2)}) {$\theta$};
  \node[cian,font=\small,anchor=south west,inner sep=6pt] at ({\R*cos(90-\th+8)},{\R*sin(90-\th+8)}) {$\dd S=R^{2}\sin\theta\,\dd\theta\,\dd\varphi$};
  % CM
  \node[punto,fill=violeta,inner sep=1.8pt] at (0,{\R/2}) {};
  \node[violeta,font=\small,anchor=west,inner sep=4pt] at (0,{\R/2}) {CM};
  \draw[violeta,line width=0.5pt,<->] (-0.25,0) -- (-0.25,{\R/2});
  \draw[aux] (-0.33,{\R/2}) -- (0,{\R/2});
  \node[violeta,font=\small,anchor=east,inner sep=3pt] at (-0.25,{\R/4}) {$R/2$};
  \node[punto,fill=tenue,inner sep=0.9pt] at (0,0) {};
  \node[peq,anchor=north,inner sep=6pt] at (0,0) {$O$};
\end{tikzpicture}
\caption{Cascarón hemisférico. Elemento de superficie en esféricas con $r=R$;
el CM está en el eje, a $R/2$ de la base.}
\label{fig:cascaron}
\end{figure}
% ---------------------------------------------------------------------

\begin{clave}[violeta]{Consejo: comprobar con casos límite}
Si la densidad depende de un parámetro, la fórmula final del CM tiene que dar
resultados obvios en los casos extremos de ese parámetro; si no los da, hay
un error. Además, el CM es un promedio de posiciones: tiene que quedar dentro
de la envolvente convexa de la masa. Sólo puede quedar fuera del propio objeto
si este es curvo, como el $2R/\pi$ del alambre semicircular (\cref{ej:semicirculo}).
\end{clave}

\begin{ejemplo}[Varilla con $\lambda=\lambda_{0}(x/L)^{n}$]
\label{ej:limites}
Para $x\in[0,L]$:
\[
  X=\frac{\int_{0}^{L}x\,\lambda\,\dd x}{\int_{0}^{L}\lambda\,\dd x}
  =\frac{L^{n+2}/(n+2)}{L^{n+1}/(n+1)}=\frac{n+1}{n+2}\,L .
\]
\begin{itemize}
  \item \emph{Densidad uniforme} ($n=0$): $X=L/2$, el punto medio.
  \item \emph{Masa concentrada en $x=L$} ($n\to\infty$): $(x/L)^{n}\to0$ salvo cerca de $x=L$, y $X\to L$.
  \item $X$ crece con $n$ y queda siempre en $[L/2,L)\subset[0,L]$ (\cref{fig:limites}).
\end{itemize}
El \cref{prob:3} pasa la misma prueba: $k=0$ da $L/2$ y $\lambda_{0}=0$ da $2L/3$, que es el
caso $n=1$ de este ejemplo.
\end{ejemplo}

% ---------------------------------------------------------------------
\begin{figure}[ht]
\centering
\begin{tikzpicture}
\begin{axis}[oscuro,width=10cm,height=5.8cm,
  xmin=0,xmax=1.05,ymin=0,ymax=1.5,xlabel={$x/L$},ylabel={$\lambda/\lambda_{0}$},
  xtick={0,0.5,0.667,0.8,0.917,1},xticklabels={$0$,$\frac12$,$\frac23$,$\frac45$,$\frac{11}{12}$,$1$},
  ytick={0,0.5,1},grid=none,legend pos=north west,legend columns=2]
  \addplot[cian,domain=0:1,samples=2] {1};
  \addlegendentry{$n=0$}
  \addplot[ambar,domain=0:1,samples=60] {x};
  \addlegendentry{$n=1$}
  \addplot[rosa,domain=0:1,samples=120] {x^3};
  \addlegendentry{$n=3$}
  \addplot[violeta,domain=0:1,samples=200] {x^10};
  \addlegendentry{$n=10$}
  \draw[cian,line width=0.6pt,dash pattern=on 2pt off 2pt] (axis cs:0.5,0) -- (axis cs:0.5,1);
  \draw[ambar,line width=0.6pt,dash pattern=on 2pt off 2pt] (axis cs:0.6667,0) -- (axis cs:0.6667,0.6667);
  \draw[rosa,line width=0.6pt,dash pattern=on 2pt off 2pt] (axis cs:0.8,0) -- (axis cs:0.8,0.512);
  \draw[violeta,line width=0.6pt,dash pattern=on 2pt off 2pt] (axis cs:0.9167,0) -- (axis cs:0.9167,0.419);
  \foreach \xx/\cc in {0.5/cian,0.6667/ambar,0.8/rosa,0.9167/violeta}{
    \edef\temp{\noexpand\node[punto,fill=\cc,inner sep=1.5pt] at (axis cs:\xx,0) {};}\temp}
\end{axis}
\end{tikzpicture}
\caption{Densidades $\lambda=\lambda_{0}(x/L)^{n}$ y su CM, $X=\frac{n+1}{n+2}L$ (líneas a trazos).
Al crecer $n$, la masa se acumula en $x=L$ y el CM se desplaza hacia allí, sin salir nunca de la varilla.}
\label{fig:limites}
\end{figure}
% ---------------------------------------------------------------------

% =====================================================================
\subsection{Tipo B · Momento angular y König}
% =====================================================================

\begin{problema}{2}[König para tres partículas]
\label{prob:2}
Para el sistema de tres partículas (de igual masa $m$) con vectores de posición
$\vb{r}_{1}=4\,\ii+(2t^{3}-4)\,\jj$,\quad $\vb{r}_{2}=\ii-2t\,\jj+6t^{3}\,\kk$,\quad
$\vb{r}_{3}=t^{2}\,\ii+4t\,\jj+3\,\kk$,
hallar las componentes del CM y relativa al mismo del vector momento angular.
\end{problema}

\begin{resolucion}
Se pide la descomposición de König \eqref{eq:konig1},
$\vb{L}_{O}=\vb{L}_{\text{orb}}+\vb{L}'$.
\begin{enumerate}
  \item \emph{Velocidades}:
  $\vb{v}_{1}=(0,\,6t^{2},\,0)$,\quad $\vb{v}_{2}=(0,\,-2,\,18t^{2})$,\quad $\vb{v}_{3}=(2t,\,4,\,0)$.
  \item \emph{CM} ($M=3m$, media de las posiciones):
  \[
    \Rc=\Bigl(\frac{t^{2}+5}{3},\ \frac{2t^{3}+2t-4}{3},\ 2t^{3}+1\Bigr),\qquad
    \Vc=\Bigl(\frac{2t}{3},\ 2t^{2}+\frac23,\ 6t^{2}\Bigr).
  \]
  \item \emph{Momento angular total}, partícula a partícula:
  \begin{align*}
    \vb{r}_{1}\times\vb{v}_{1}&=(0,\,0,\,24t^{2}),\\
    \vb{r}_{2}\times\vb{v}_{2}&=(-24t^{3},\,-18t^{2},\,-2),\\
    \vb{r}_{3}\times\vb{v}_{3}&=(-12,\,6t,\,-4t^{2}),
  \end{align*}
  \[
    \vb{L}_{O}=m\sum_{\alpha}\vb{r}_{\alpha}\times\vb{v}_{\alpha}
    =m\,\bigl(-24t^{3}-12,\ -18t^{2}+6t,\ 20t^{2}-2\bigr).
  \]
  \item \emph{Parte orbital} (del CM): $\vb{L}_{\text{orb}}=\Rc\times M\Vc=3m\,\Rc\times\Vc$,
  \[
    \resultado{\vb{L}_{\text{orb}}=m\Bigl(8t^{3}-30t^{2}-2,\ \ -2t^{4}-30t^{2}+2t,\ \ \tfrac13(2t^{4}+28t^{2}+8t+10)\Bigr)} .
  \]
  \item \emph{Parte relativa al CM}, por diferencia (König):
  \[
    \resultado{\vb{L}'=\vb{L}_{O}-\vb{L}_{\text{orb}}
    =m\Bigl(-32t^{3}+30t^{2}-10,\ \ 2t^{4}+12t^{2}+4t,\ \ \tfrac13(-2t^{4}+32t^{2}-8t-16)\Bigr)} .
  \]
  \item \emph{Comprobación}: calculando directamente
  $\vb{L}'=\sum m\,\rp_{\alpha}\times\vp_{\alpha}$ con $\rp_{\alpha}=\vb{r}_{\alpha}-\Rc$,
  $\vp_{\alpha}=\vb{v}_{\alpha}-\Vc$, se obtiene lo mismo. Por ejemplo, en $t=0$:
  $\vb{L}_{O}=m(-12,0,-2)$, $\vb{L}_{\text{orb}}=m(-2,0,\tfrac{10}{3})$ y $\vb{L}'=m(-10,0,-\tfrac{16}{3})$.\qedhere
\end{enumerate}
\end{resolucion}

\begin{clave}[violeta]{Consejo}
Calcula siempre $\vb{L}_{O}$ y $\vb{L}_{\text{orb}}$, que son fáciles, y obtén $\vb{L}'$ por
diferencia; usa el cálculo directo de $\vb{L}'$ sólo como comprobación en un
instante concreto. Por el \cref{teo:momento-angular}, $\dot{\vb{L}}'$ da el momento de
las fuerzas externas respecto del CM.
\end{clave}

% =====================================================================
\subsection{Tipo C · Conservación del momento y movimiento del CM}
% =====================================================================

\begin{problema}{5}[La granada]
\label{prob:5}
Una granada es lanzada al aire y explota en tres pedazos. Dos de ellos, con
masas iguales a la cuarta parte de la masa de la granada, salen disparados
con direcciones perpendiculares entre sí tras la explosión, siendo la
velocidad final de cada uno, medida respecto al centro de masas, $80$\,m/s.
Determinar la velocidad del tercer pedazo después de la explosión, suponiendo
despreciable la masa del explosivo.
\end{problema}

\begin{resolucion}
\begin{enumerate}
  \item \emph{Masas}: $m_{1}=m_{2}=M/4$, luego $m_{3}=M-M/2=M/2$.
  \item \emph{Momento}: la explosión es interna y dura muy poco, así que $\vb{P}$
  se conserva (el peso no cambia $\vb{P}$ durante el instante de la explosión).
  En el sistema del CM, $\vb{P}'=\vb{0}$ antes y después (\cref{cor:conserva-P}):
  \[
    \frac M4\,\vp_{1}+\frac M4\,\vp_{2}+\frac M2\,\vp_{3}=\vb{0}
    \qquad\Longrightarrow\qquad
    \vp_{3}=-\frac{\vp_{1}+\vp_{2}}{2}.
  \]
  \item \emph{Ejes} (\cref{fig:granada}): $\vp_{1}=80\,\ii$ y $\vp_{2}=80\,\jj$ (m/s), que son
  perpendiculares:
  \[
    \resultado{\vp_{3}=-40\,(\ii+\jj)\ \text{m/s},\qquad
    v'_{3}=40\sqrt2\simeq56{,}6\ \text{m/s}} ,
  \]
  dirigida según la bisectriz opuesta, a $135^{\circ}$ de cada uno de los otros dos.
  \item \emph{En el laboratorio}, si el CM llevaba en ese instante la
  velocidad $\Vc$ de la granada, $\vb{v}_{3}=\Vc+\vp_{3}$. El CM sigue la misma
  parábola que habría seguido la granada sin explotar.\qedhere
\end{enumerate}
\end{resolucion}

% ---------------------------------------------------------------------
\begin{figure}[ht]
\centering
\begin{tikzpicture}[scale=1.15]
  \coordinate (C) at (0,0);
  \draw[vec,cian] (C) -- (2.0,0) node[anchor=west,etiq,inner sep=3pt]{$\vp_{1}$ \ {\color{tenue}\footnotesize($M/4$)}};
  \draw[vec,rosa] (C) -- (0,2.0) node[anchor=south,etiq,inner sep=3pt]{$\vp_{2}$ \ {\color{tenue}\footnotesize($M/4$)}};
  \draw[vec,ambar] (C) -- (-1.0,-1.0) node[anchor=north east,etiq,inner sep=2pt]{$\vp_{3}$ \ {\color{tenue}\footnotesize($M/2$)}};
  \draw[aux] (C) -- (1.4,1.4);
  \draw[texto!60,line width=0.4pt] (0.25,0) -- (0.25,0.25) -- (0,0.25);
  \node[peq,anchor=south west,inner sep=1pt] at (1.4,1.4) {$\vp_{1}+\vp_{2}$};
  \node[punto,fill=violeta,inner sep=1.8pt] at (C) {};
  \node[violeta,font=\small,anchor=north west,inner sep=3pt] at (C) {CM};
\end{tikzpicture}
\caption{Explosión vista desde el CM: $\vb{P}'=\vb{0}$ obliga a que el tercer
fragmento, de masa doble, salga en sentido opuesto a $\vp_{1}+\vp_{2}$ con la mitad de
su módulo.}
\label{fig:granada}
\end{figure}
% ---------------------------------------------------------------------

\begin{problema}{6}[Máquina de Atwood: el CM]
\label{prob:6}
Dos bloques de masas $m_{1}$ y $m_{2}$ están unidos mediante un hilo
inextensible y sin masa que pasa por la garganta de una polea lisa. Hallar la
velocidad y aceleración del centro de masas cuando se abandona el sistema
partiendo del reposo.
\end{problema}

\begin{resolucion}
Supongamos $m_{1}>m_{2}$; $m_{1}$ baja y $m_{2}$ sube (\cref{fig:atwood}). El eje $z$ es
vertical hacia arriba y $a$ es el módulo de la aceleración de cada bloque.
\begin{enumerate}
  \item \emph{Bloques} (polea lisa y sin masa: misma tensión $T$ a ambos lados):
  \[
    m_{1}g-T=m_{1}a,\qquad T-m_{2}g=m_{2}a
    \qquad\Longrightarrow\qquad
    a=\frac{m_{1}-m_{2}}{m_{1}+m_{2}}\,g,\qquad T=\frac{2m_{1}m_{2}}{m_{1}+m_{2}}\,g .
  \]
  \item \emph{CM}: $Z=\dfrac{m_{1}z_{1}+m_{2}z_{2}}{m_{1}+m_{2}}$, con $\ddot z_{1}=-a$ y $\ddot z_{2}=+a$ (hilo inextensible):
  \[
    \ddot Z=\frac{-m_{1}a+m_{2}a}{m_{1}+m_{2}}
    =\resultado{-\Bigl(\frac{m_{1}-m_{2}}{m_{1}+m_{2}}\Bigr)^{2}g},
    \qquad
    \resultado{\dot Z(t)=-\Bigl(\frac{m_{1}-m_{2}}{m_{1}+m_{2}}\Bigr)^{2}g\,t},
  \]
  dirigidas hacia abajo, partiendo del reposo. La componente horizontal del CM es fija.
  \item \emph{Comprobación con el \cref{teo:momento-lineal}}. Sistema: bloques $+$ hilo.
  Las fuerzas externas son los pesos, $-(m_{1}+m_{2})g\,\kk$, y la de la polea sobre el
  hilo, $+2T\,\kk$:
  \[
    M\ddot Z=2T-Mg=\frac{4m_{1}m_{2}-(m_{1}+m_{2})^{2}}{m_{1}+m_{2}}\,g=-\frac{(m_{1}-m_{2})^{2}}{m_{1}+m_{2}}\,g .\ \checkmark
  \]
  \item \emph{Lectura}: el CM cae con aceleración menor que $g$; si $m_{1}=m_{2}$, el CM
  queda en reposo.\qedhere
\end{enumerate}
\end{resolucion}

% ---------------------------------------------------------------------
\begin{figure}[ht]
\centering
\begin{tikzpicture}[scale=1.25]
  \def\r{0.75}
  % soporte y polea
  \draw[texto!60,line width=0.8pt] (-1.5,4.2) -- (1.5,4.2);
  \foreach \x in {-1.4,-1.1,...,1.5}{\draw[texto!40,line width=0.4pt] (\x,4.2) -- (\x+0.2,4.4);}
  \draw[texto!60,line width=0.6pt] (0,4.2) -- (0,3.3);
  \draw[texto!85,line width=0.9pt,fill=caja] (0,3.3) circle (\r);
  \node[punto,fill=tenue,inner sep=1pt] at (0,3.3) {};
  % hilos
  \draw[texto!85,line width=0.8pt] (-\r,3.3) -- (-\r,0.7);
  \draw[texto!85,line width=0.8pt] (\r,3.3) -- (\r,2.2);
  % bloques
  \draw[texto!85,line width=0.8pt,fill=caja] (-1.25,-0.3) rectangle (-0.25,0.7);
  \node[etiq] at (-\r,0.2) {$m_{1}$};
  \draw[texto!85,line width=0.8pt,fill=caja] (0.35,1.5) rectangle (1.15,2.2);
  \node[etiq] at (\r,1.85) {$m_{2}$};
  % tensiones sobre los bloques (en el hilo)
  \draw[vec,cian] (-\r,0.7) -- (-\r,1.55);
  \node[cian,font=\small,anchor=east,inner sep=3pt] at (-\r,1.25) {$T$};
  \draw[vec,cian] (\r,2.2) -- (\r,2.85);
  \node[cian,font=\small,anchor=west,inner sep=3pt] at (\r,2.6) {$T$};
  % aceleraciones
  \draw[vec,rosa] (-1.65,0.4) -- (-1.65,-0.4) node[anchor=north,etiq,inner sep=2pt]{$a$};
  \draw[vec,rosa] (1.55,1.6) -- (1.55,2.4) node[anchor=south,etiq,inner sep=2pt]{$a$};
  % CM (m1=3, m2=1)
  \pgfmathsetmacro{\zc}{(0.2*3+1.85*1)/4}
  \node[punto,fill=violeta,inner sep=1.7pt] at (0,\zc) {};
  \draw[vec,violeta] (0,\zc) -- (0,{\zc-0.6});
  \node[violeta,font=\footnotesize,anchor=west,inner sep=3pt] at (0,{\zc-0.3}) {$\ddot Z$};
  \node[violeta,font=\footnotesize,anchor=south west,inner sep=2pt] at (0,\zc) {CM};
\end{tikzpicture}
\caption{Máquina de Atwood con $m_{1}>m_{2}$: los bloques tienen aceleraciones
opuestas de módulo $a$. El CM (violeta) cae con
$\abs{\ddot Z}=\bigl(\tfrac{m_{1}-m_{2}}{m_{1}+m_{2}}\bigr)^{2}g<g$.}
\label{fig:atwood}
\end{figure}
% ---------------------------------------------------------------------

\begin{problema}{7}[Colapso gravitatorio de dos cuerpos]
\label{prob:7}
Dos partículas de masas $m_{1}$ y $m_{2}$ están inicialmente en reposo separadas una
distancia $D$ y sometidas tan sólo a su acción gravitatoria mutua. Calcular el
tiempo que transcurrirá hasta que colisionen.
\end{problema}

\begin{resolucion}
\begin{enumerate}
  \item \emph{CM}: no hay fuerzas externas y parten del reposo, así que el CM
  está quieto (\cref{cor:conserva-P}). Chocan en el CM.
  \item \emph{Ecuación relativa}. Con $\vb{r}\defeq\vb{r}_{1}-\vb{r}_{2}$ y $r=\abs{\vb{r}}$:
  \[
    m_{1}\ddot{\vb{r}}_{1}=-\frac{Gm_{1}m_{2}}{r^{2}}\,\uv{r},\qquad
    m_{2}\ddot{\vb{r}}_{2}=+\frac{Gm_{1}m_{2}}{r^{2}}\,\uv{r}
    \qquad\Longrightarrow\qquad
    \ddot{\vb{r}}=-\frac{G(m_{1}+m_{2})}{r^{2}}\,\uv{r},
  \]
  restando tras dividir cada ecuación por su masa. Equivale a una partícula
  atraída por una masa fija $M\defeq m_{1}+m_{2}$.
  \item \emph{Energía} (movimiento radial, $\dot r(0)=0$, $r(0)=D$):
  \[
    \tfrac12\dot r^{2}-\frac{GM}{r}=-\frac{GM}{D}
    \qquad\Longrightarrow\qquad
    \dot r=-\sqrt{2GM\Bigl(\frac1r-\frac1D\Bigr)}=-\sqrt{\frac{2GM}{D}}\sqrt{\frac{D-r}{r}} .
  \]
  \item \emph{Cuadratura} \eqref{eq:cuadratura}:
  \[
    t_{c}=\sqrt{\frac{D}{2GM}}\int_{0}^{D}\sqrt{\frac{r}{D-r}}\;\dd r .
  \]
  \item \emph{Cambio} $r=D\sin^{2}\theta$, con $\dd r=2D\sin\theta\cos\theta\,\dd\theta$ y
  $\sqrt{r/(D-r)}=\tan\theta$:
  \[
    \int_{0}^{D}\sqrt{\frac{r}{D-r}}\,\dd r=\int_{0}^{\pi/2}2D\sin^{2}\theta\,\dd\theta=\frac{\pi D}{2}
    \qquad\Longrightarrow\qquad
    \resultado{t_{c}=\frac{\pi}{2}\sqrt{\frac{D^{3}}{2G(m_{1}+m_{2})}}} .
  \]
  \item \emph{Comprobación con Kepler}: la caída radial es una elipse degenerada
  de semieje mayor $a=D/2$, y $t_{c}$ es medio periodo:
  $\tfrac12\cdot2\pi\sqrt{a^{3}/GM}=\pi\sqrt{D^{3}/(8GM)}$, que coincide.\qedhere
\end{enumerate}
\end{resolucion}

\begin{clave}[violeta]{Consejo}
En problemas de dos cuerpos aislados: el CM se mueve uniformemente (o está
quieto) y la coordenada relativa $\vb{r}=\vb{r}_{1}-\vb{r}_{2}$ obedece la ecuación de
una sola partícula. La energía da $\dot r(r)$ y la integral \eqref{eq:cuadratura}, el
tiempo.
\end{clave}

% =====================================================================
\subsection{Tipo D · Masa variable}
% =====================================================================

\begin{clave}{Método: masa variable}
\begin{enumerate}[leftmargin=1.4em]
  \item Identificar $\dot m$ (signo) y la velocidad $\vb{u}$ de la masa intercambiada
        \emph{relativa al cuerpo}.
  \item Escribir $m\dot{\vb{v}}=\vb{F}\ext+\dot m\,\vb{u}$ y proyectar sobre ejes.
  \item Si se capta masa en reposo, usar $\od{(m\vb{v})}{t}=\vb{F}\ext$ \eqref{eq:captacion}.
\end{enumerate}
\end{clave}

\begin{problema}{8}[Gota que crece en una nube]
\label{prob:8}
Una gota esférica de agua pasa por una nube y varía su masa de modo que su
radio crece como $r=r_{0}(1+bt)$, $b>0$. (a) Demostrar que $\dd m/\dd t$ es
proporcional a su superficie. (b) Despreciando la resistencia del aire, obtener
su velocidad en función del tiempo si parte del reposo.
\end{problema}

\begin{resolucion}
\begin{enumerate}
  \item \emph{(a)} Con $m=\tfrac43\pi\rho_{w}r^{3}$ y $\dot r=r_{0}b$ constante:
  \[
    \od{m}{t}=4\pi\rho_{w}r^{2}\,\dot r=\resultado{(\rho_{w}r_{0}b)\,S},\qquad S=4\pi r^{2}.
  \]
  La gota capta agua a un ritmo proporcional a su superficie de captura.
  \item \emph{(b) Ecuación}. Las gotitas de la nube están en reposo, luego por
  \eqref{eq:captacion}, con $\vb{F}\ext=m\vb{g}$ y la vertical hacia abajo:
  \[
    \od{(mv)}{t}=mg .
  \]
  \item \emph{Integrar}. Con $m=m_{0}(1+bt)^{3}$ y $v(0)=0$:
  \[
    m_{0}(1+bt)^{3}v=m_{0}g\int_{0}^{t}(1+bt')^{3}\dd t'=m_{0}g\,\frac{(1+bt)^{4}-1}{4b}.
  \]
  \item \emph{Resultado}:
  \[
    \resultado{v(t)=\frac{g}{4b}\Bigl[(1+bt)-\frac{1}{(1+bt)^{3}}\Bigr]} .
  \]
  \item \emph{Comprobaciones}: para $bt\ll1$, $v\simeq gt-\tfrac32bg\,t^{2}$ (caída libre al
  principio); para $bt\gg1$, $v\simeq\tfrac{g}{4b}(1+bt)$, con aceleración límite $g/4$.
  La masa captada, que hay que acelerar desde el reposo, frena la gota.\qedhere
\end{enumerate}
\end{resolucion}

% ---------------------------------------------------------------------
\begin{figure}[ht]
\centering
\begin{tikzpicture}
\begin{axis}[oscuro,width=9.5cm,height=6cm,
  xmin=0,xmax=6,ymin=0,ymax=3.2,xlabel={$bt$},ylabel={$v\,b/g$},
  legend pos=south east]
  \addplot[cian,domain=0:3.2,samples=80] {x};
  \addlegendentry{caída libre, $v=gt$}
  \addplot[ambar,domain=0:6,samples=150] {((1+x)-1/(1+x)^3)/4};
  \addlegendentry{gota que crece}
  \addplot[rosa,dashed,domain=0:6,samples=10] {(1+x)/4};
  \addlegendentry{asíntota, pendiente $g/4$}
\end{axis}
\end{tikzpicture}
\caption{Velocidad de la gota (\cref{prob:8}). Al principio cae libremente; al
crecer, la masa captada en reposo la frena y la aceleración tiende a $g/4$.}
\label{fig:gota}
\end{figure}
% ---------------------------------------------------------------------

\begin{problema}{9}[Cohete que desciende]
\label{prob:9}
Un cohete autopropulsado se acerca a la superficie terrestre formando un
ángulo de $30^{\circ}$ con la horizontal y encontrando en el aire una resistencia
igual a la tercera parte de su peso. Calcular su aceleración si expulsa gases
a razón de $G$\,kg/s con velocidad $v_{c}$ respecto al cohete.
\end{problema}

\begin{resolucion}
Sea $m$ la masa instantánea, $\utau$ la dirección del movimiento (hacia abajo,
a $30^{\circ}$ de la horizontal) y $\un$ la normal en el plano vertical
(\cref{fig:cohete30}). Supuestos: el gas se expulsa hacia atrás,
$\vb{u}=-v_{c}\,\utau$, con $\dot m=-G$; la resistencia es $-\tfrac13mg\,\utau$.
\begin{enumerate}
  \item \emph{Ecuación} \eqref{eq:masa-variable}:
  \[
    m\vb{a}=m\vb{g}-\tfrac13mg\,\utau+(-G)(-v_{c}\,\utau)
    =m\vb{g}+\Bigl(Gv_{c}-\tfrac13mg\Bigr)\utau .
  \]
  \item \emph{Peso descompuesto}: $\vb{g}=g\sin30^{\circ}\,\utau+g\cos30^{\circ}\,\un=\tfrac g2\,\utau+\tfrac{\sqrt3}{2}g\,\un$.
  \item \emph{Componentes}:
  \[
    \resultado{a_{T}=\frac g2-\frac g3+\frac{Gv_{c}}{m}=\frac g6+\frac{Gv_{c}}{m}},\qquad
    \resultado{a_{N}=\frac{\sqrt3}{2}\,g},
  \]
  y el módulo es $a=\sqrt{a_{T}^{2}+\tfrac34g^{2}}$.
  \item \emph{Variante}: si el cohete frena expulsando el gas hacia delante
  ($\vb{u}=+v_{c}\,\utau$), el empuje cambia de signo y $a_{T}=\tfrac g6-\tfrac{Gv_{c}}{m}$. Para
  que el cohete mantenga una trayectoria rectilínea, otra fuerza (sustentación
  o empuje lateral) tiene que anular $a_{N}$.\qedhere
\end{enumerate}
\end{resolucion}

% ---------------------------------------------------------------------
\begin{figure}[ht]
\centering
\begin{tikzpicture}[scale=1.25]
  \coordinate (P) at (2.2,2.6);
  % trayectoria y horizontal
  \draw[aux] ($(P)+(150:1.9)$) -- ($(P)+(-30:3.4)$);
  \draw[aux] (P) -- ($(P)+(2.6,0)$);
  \draw[violeta,line width=0.6pt] ($(P)+(2.2,0)$) arc[start angle=0,end angle=-30,radius=2.2];
  \node[violeta,font=\small,anchor=west,inner sep=3pt] at ($(P)+(-15:2.2)$) {$30^{\circ}$};
  % cohete
  \begin{scope}[shift={(P)},rotate=-30]
    \draw[texto!85,line width=0.8pt,fill=caja] (-0.8,-0.2) -- (0.6,-0.2) -- (0.9,0) -- (0.6,0.2) -- (-0.8,0.2) -- cycle;
    \foreach \p in {(-1.05,0.08),(-1.25,-0.1),(-1.45,0.12),(-1.6,-0.05)} \fill[naranja!80] \p circle (0.055);
  \end{scope}
  \node[punto,inner sep=1.2pt] at (P) {};
  % fuerzas y vectores
  \draw[vec,cian] ($(P)+(-30:0.95)$) -- ++(-30:1.0) node[anchor=north east,etiq,inner sep=1pt]{$\utau$};
  \draw[vec,rosa] (P) -- ++(0,-1.4) node[anchor=north,etiq,inner sep=2pt]{$m\vb{g}$};
  \draw[vec,ambar] ($(P)+(60:0.55)$) -- ($(P)+(60:0.55)+(150:1.1)$) node[anchor=south,etiq,inner sep=3pt]{$\tfrac13 mg$};
  \draw[vec,verde] (P) -- ++(-120:1.0) node[anchor=north east,etiq,inner sep=1pt]{$\un$};
  % suelo
  \draw[texto!50,line width=0.6pt] (0,0) -- (7.2,0);
  \foreach \x in {0.2,0.6,...,7.1}{\draw[texto!35,line width=0.4pt] (\x,0) -- (\x-0.2,-0.2);}
\end{tikzpicture}
\caption{Cohete que desciende a $30^{\circ}$: peso, resistencia opuesta al
movimiento y empuje del gas expulsado hacia atrás (a lo largo de $\utau$).}
\label{fig:cohete30}
\end{figure}
% ---------------------------------------------------------------------

% =====================================================================
\subsection{Tipo E · Virial y Pappus}
% =====================================================================

\begin{problema}{10}[Virial para osciladores]
\label{prob:10}
Verificar el cumplimiento del teorema del virial para el caso del oscilador
armónico y anarmónico (con término cuártico) unidimensional.
\end{problema}

\begin{resolucion}
\begin{enumerate}
  \item \emph{Virial en 1D}: $\prom{T}=\tfrac12\prom{x\,V'(x)}$, por \eqref{eq:virial} con $F=-V'$.
  \item \emph{Armónico}, $V=\tfrac12kx^{2}$: $xV'=kx^{2}=2V$, así que el virial predice
  $\prom{T}=\prom{V}=E/2$. \emph{Verificación explícita} con $x=A\cos\omega t$:
  \begin{gather*}
    T=\tfrac12kA^{2}\sin^{2}\omega t,\qquad V=\tfrac12kA^{2}\cos^{2}\omega t,\qquad
    \prom{\sin^{2}}=\prom{\cos^{2}}=\tfrac12\\
    \Longrightarrow\ \resultado{\prom{T}=\prom{V}=\tfrac14kA^{2}=\tfrac E2}.
  \end{gather*}
  \item \emph{Anarmónico}, $V=\tfrac12kx^{2}+\tfrac14\lambda x^{4}$:
  $xV'=kx^{2}+\lambda x^{4}=2V_{2}+4V_{4}$, donde $V_{2}=\tfrac12kx^{2}$ y $V_{4}=\tfrac14\lambda x^{4}$. Luego
  \[
    \resultado{\prom{T}=\prom{V_{2}}+2\prom{V_{4}}} .
  \]
  \item \emph{Caso cuártico puro} ($k=0$): $n=4$ en \eqref{eq:virial-n}, y
  $\prom{T}=2\prom{V}$, es decir, $\prom{T}=\tfrac23E$ y $\prom{V}=\tfrac13E$.
  \item \emph{Verificación numérica} (\cref{fig:virial}): integrando la ecuación
  de movimiento durante dos periodos exactos se obtiene
  $\prom{T}/E=0{,}6667$ y $\prom{V}/E=0{,}3333$. El periodo del cuártico depende de la amplitud
  (no es isócrono), pero la proporción entre $\prom{T}$ y $\prom{V}$ no.\qedhere
\end{enumerate}
\end{resolucion}

% ---------------------------------------------------------------------
\begin{figure}[ht]
\centering
\begin{tikzpicture}
\begin{groupplot}[group style={group size=2 by 1,horizontal sep=12mm},
  oscuro,width=7.2cm,height=5.4cm,xmin=0,xmax=2,ymin=0,ymax=1.08,
  xlabel={$t/T_{0}$},xtick={0,0.5,1,1.5,2},height=5.8cm]
  \nextgroupplot[title={(a) Armónico: $\prom{T}=\prom{V}=\frac E2$},ylabel={energía$/E$},
    ytick={0,0.5,1},yticklabels={$0$,$\frac12$,$1$}]
    \addplot[cian] table[x=t,y=T]{datos/armonico.dat};
    \addplot[rosa] table[x=t,y=V]{datos/armonico.dat};
    \addplot[texto!70,dashed,line width=0.7pt,domain=0:2,samples=2] {0.5};
  \nextgroupplot[title={(b) Cuártico: $\prom{T}=\frac{2E}{3}$, $\prom{V}=\frac E3$},
    ytick={0,0.3333,0.6667,1},yticklabels={$0$,$\frac13$,$\frac23$,$1$}]
    \addplot[cian] table[x=t,y=T]{datos/cuartico.dat};
    \addplot[rosa] table[x=t,y=V]{datos/cuartico.dat};
    \addplot[cian!70,dashed,line width=0.7pt,domain=0:2,samples=2] {0.6667};
    \addplot[rosa!70,dashed,line width=0.7pt,domain=0:2,samples=2] {0.3333};
\end{groupplot}
\end{tikzpicture}
\caption{Energía cinética (cian) y potencial (rosa) durante dos periodos,
normalizadas por $E$. Las rectas a trazos son los promedios: iguales en el
armónico ($n=2$) y en proporción $2:1$ en el cuártico ($n=4$), como predice
$2\prom{T}=n\prom{V}$.}
\label{fig:virial}
\end{figure}
% ---------------------------------------------------------------------

\begin{problema}{11}[Volumen del cono por Pappus]
\label{prob:11}
Usando el teorema de Pappus, calcular el volumen de un cono homogéneo de
base circular de radio $R$ y altura $h$.
\end{problema}

\begin{resolucion}
\begin{enumerate}
  \item \emph{Región generatriz}: el triángulo rectángulo de catetos $R$ (base,
  perpendicular al eje) y $h$ (sobre el eje) genera el cono al girar alrededor
  del cateto $h$ (\cref{fig:cono}). Su área es $A=\tfrac12Rh$.
  \item \emph{Centroide}: el de un triángulo es la media de sus vértices
  (\cref{prob:1}). Con vértices $(0,0)$, $(R,0)$ y $(0,h)$, siendo la primera coordenada la
  distancia al eje,
  \[
    \bar\rho=\frac{0+R+0}{3}=\frac R3 .
  \]
  \item \emph{Pappus II} (\cref{teo:pappus2}):
  \[
    V=2\pi\bar\rho\,A=2\pi\cdot\frac R3\cdot\frac{Rh}{2}=\resultado{\frac13\pi R^{2}h} .
  \]
  \item \emph{Comprobación} por integración (discos):
  $V=\int_{0}^{h}\pi\bigl(R(1-z/h)\bigr)^{2}\dd z=\pi R^{2}h\int_{0}^{1}(1-u)^{2}\dd u=\tfrac13\pi R^{2}h$.\qedhere
\end{enumerate}
\end{resolucion}

% ---------------------------------------------------------------------
\begin{figure}[ht]
\centering
\begin{tikzpicture}[scale=1.2]
  \def\R{1.8}\def\h{2.8}
  \draw[texto!45,line width=0.6pt] (0,0) ellipse[x radius=\R,y radius=0.36];
  \draw[texto!45,line width=0.6pt] (-\R,0) -- (0,\h) -- (\R,0);
  \fill[cian,opacity=0.18] (0,0) -- (\R,0) -- (0,\h) -- cycle;
  \draw[cian,line width=1pt] (0,0) -- (\R,0) -- (0,\h) -- cycle;
  \draw[eje] (0,-0.6) -- (0,\h+0.6) node[anchor=south,peq,inner sep=2pt]{eje};
  \coordinate (G) at ({\R/3},{\h/3});
  \node[punto,fill=violeta,inner sep=1.7pt] at (G) {};
  \node[violeta,font=\small,anchor=south west,inner sep=2pt] at (G) {$G$};
  \draw[violeta,line width=0.6pt,dash pattern=on 3pt off 2pt] (0,{\h/3}) ellipse[x radius={\R/3},y radius=0.12];
  \draw[violeta,line width=0.5pt,<->] (0,-0.62) -- ({\R/3},-0.62);
  \draw[aux,violeta] ({\R/3},-0.7) -- (G);
  \node[violeta,font=\footnotesize,anchor=north,inner sep=2pt] at ({\R/6},-0.62) {$R/3$};
  \node[peq,anchor=north west,inner sep=2pt] at ({\R/2+0.4},-0.3) {$R$};
  \draw[tenue,line width=0.5pt,<->] (-2.15,0) -- (-2.15,\h);
  \node[peq,anchor=east,inner sep=3pt] at (-2.15,{\h/2}) {$h$};
\end{tikzpicture}
\caption{Cono generado por el triángulo rectángulo (cian) al girar alrededor
del cateto $h$. Su centroide $G$ está a $R/3$ del eje y recorre $2\pi R/3$.}
\label{fig:cono}
\end{figure}
% ---------------------------------------------------------------------

\begin{clave}[violeta]{Consejo}
Pappus funciona en los dos sentidos: si conoces el volumen (o el área)
generado, despejas la posición del centroide sin integrar
(\cref{ej:pappus-centroides}).
\end{clave}
